\subsection{OPEN EQUIVALENCE RELATIONS AND CLOSED EQUIVALENCE RELATIONS}

DEFINITION 2. An equivalence relation R on a topological space X is said to be open (resp. closed) if the canonical mapping of X onto X/R is open (resp. closed).

It comes to the same thing to say that the saturation of each open (resp. closed) subset of X with respect to R is open (resp. closed) in X (§ 3, no. 4).

Examples. 1) Let X be a topological space, \(\Gamma\) a group of homeomorphisms of X onto itself, and let R be the equivalence relation

\[
" \text { there   exists } \sigma \in \Gamma \text { such   that } y = \sigma (x)"
\]

between x and y (thus R is the equivalence relation whose classes are the orbits of \(\Gamma\) in X. The relation R is open, for the saturation of a subset A of X with respect to R is open so is each \(\sigma(\mathbf{A})\) and hence so is their union.

* On the real line R the equivalence relation \(x \equiv y \pmod{1}\) is open, since it is derived as above from the group of translations

\[
x \rightarrow x + n \quad (n \in \mathbf {Z})
\]

(see Chapter III, § 2, no. 4). *

\begin{enumerate}
\def\labelenumi{\arabic{enumi})}
\setcounter{enumi}{1}
\tightlist
\item
  Let X be the sum of a family \((\mathbf{X}_{t})\) of subspaces of X, and let X/R be the space obtained by pasting together the \(X_{t}\) along open subsets \(A_{tx}\) by means of bijections \(h_{xt}\) (§ 2, no. 5); and suppose that \(h_{xt}\) is a homeomorphism of \(A_{tx}\) onto \(A_{xt}\) for each pair of indices \((t, x)\) .
\end{enumerate}

Then the relation R is open. For if U is open in X, the saturation of U is the union of the \(h_{xt}(\mathbf{U} \cap \mathbf{A}_{\mathbf{tx}})\) ; since \(U \cap A_{tx}\) is open in \(A_{tx}\) , \(h_{xt}(\mathbf{U} \cap \mathbf{A}_{\mathbf{tx}})\) is open in \(A_{xt}\) and therefore in X.

\begin{enumerate}
\def\labelenumi{\arabic{enumi})}
\setcounter{enumi}{2}
\tightlist
\item
  With the notation of Example 2, suppose now that the \(\mathbf{A}_{\mathrm{tx}}\) are closed and the \(h_{xt}\) are homeomorphisms; further suppose that for each index \(t\) there is only a finite number of indices \(x\) such that \(\mathbf{A}_{\mathrm{tx}} \neq \emptyset\) (i.e.~each \(\mathbf{X}_t\) is ``stuck'' to only a finite number of \(\mathbf{X}_x\) ). Then the relation R is closed. For if F is any closed subset of X, the saturation of F is the union of the sets \(h_{xt}(\mathbf{F} \cap \mathbf{A}_{\mathrm{tx}}) \subset \mathbf{A}_{xt}\) ; the assumptions made imply that this family is locally finite, and \(h_{xt}(\mathbf{F} \cap \mathbf{A}_{\mathrm{tx}})\) is closed in \(\mathbf{A}_{xt}\) and therefore in X. The result therefore follows from Proposition 4 of § 2, no. 5.
\end{enumerate}

PROPOSITION 3. Let X, Y be two topological spaces, let \(f: \mathbf{X} \to \mathbf{Y}\) be a continuous mapping, let R be the equivalence relation \(f(x) = f(y)\) on X, and let \(\mathbf{X} \xrightarrow{p} \mathbf{X}/\mathbf{R} \xrightarrow{h} f(\mathbf{X}) \xrightarrow{i} \mathbf{Y}\) be the canonical decomposition of \(f\) . Then the following three statements are equivalent:

\begin{enumerate}
\def\labelenumi{\alph{enumi})}
\item
  \(f\) is an open mapping.
\item
  The three mappings \(p, h, i\) are open.
\item
  The equivalence relation R is open, h is a homeomorphism, and \(f(\mathbf{X})\) is an open subset of Y.
\end{enumerate}

Also the preceding proposition remains true if ``open'' is replaced by ``closed'' throughout.

Since the injection \(i\) is continuous, it follows from Proposition 1c) of no. 1 that if \(f\) is open then so is \(h \circ p\) ; since \(p\) is surjective and continuous, Proposition 1b) shows that \(h\) is open; \(h\) is in any case a continuous bijection; thus \(h\) is a homeomorphism, and therefore, from Proposition 1a) of no. 1, \(p = \overline{h}^1 \circ (h \circ p)\) is an open mapping. Also {[}no. 1, Proposition 1b){]} \(i \circ h\) is open; hence {[}no. 1, Proposition 1a){]} so is \(i = (i \circ h) \circ \overline{h}^1\) . This proves that a) implies b). Conversely, Proposition 1a) of no. 1 shows that b) implies a). Finally, the equivalence of b) and c) follows immediately from the definitions.

The proof in the case of closed mappings is analogous, mutatis mutandis.

PROPOSITION 4. Let R be an open (resp. closed) equivalence relation on a topological space X, and f the canonical mapping X → X/R. Let A be a subset of X and suppose that one of the following two conditions is satisfied: a) A is open (resp. closed) in X.

\begin{enumerate}
\def\labelenumi{\alph{enumi})}
\setcounter{enumi}{1}
\tightlist
\item
  A is saturated with respect to R.
\end{enumerate}

Then the relation \(R_{A}\) induced on A is open (resp. closed) and the canonical mapping of \(A/R_{A}\) onto \(f(A)\) is a homeomorphism.

Consider the commutative diagram (1) of § 3, no. 6, which gives the canonical decomposition of \(f \circ j\) . Under condition a), \(j\) is open (resp. closed) and so is \(f\) by hypothesis; hence \(f \circ j\) is open (resp. closed) {[}no. 1, Proposition 1 a){]}, and the result follows from Proposition 3. Under condition b) we have

\[
\mathbf {A} = \overline {{{f}}} ^ {1} (f (\mathbf {A})),
\]

and \(h \circ \varphi\) is the mapping of A into \(f(A)\) which agrees with \(f\) on A; by virtue of Proposition 2 a) of no. 1, \(h \circ \varphi\) is open (resp. closed), and the result again follows from Proposition 3, applied to \(h \circ \varphi\) .

